\documentclass[12pt,a4paper]{article}
\usepackage[utf8]{inputenc}
\usepackage{amsmath, amssymb, amsfonts}
\usepackage{graphicx}
\usepackage{hyperref}
\usepackage{booktabs}
\usepackage{caption}
\usepackage{subcaption}
\usepackage{geometry}
\usepackage{multirow}
\usepackage{array}
\usepackage{setspace}
\geometry{top=2.5cm, bottom=2.5cm, left=3cm, right=3cm}
\onehalfspacing

\title{Geometry of Reasoning Trajectories in Recurrent-Depth Transformers:\\A Comprehensive Analysis of Latent Dynamics in Huginn-3.5B}
\author{Alexandr Shianov$^{1}$ \and David Sverdlov$^{1}$ \\
\small{$^{1}$Independent Researchers} \\
\small{Curator: S. Barannikov; Mentor: N. Kurdyukov}}
\date{}

\begin{document}

\maketitle

\begin{abstract}
We present an inference-only analysis of latent trajectories in the Huginn-3.5B recurrent-depth transformer (Geiping et al., arXiv:2502.05171). Without fine-tuning, we extract hidden states at each recurrence step via hooks on \texttt{core\_block[-1]} and compute geometric, dynamical, and topological descriptors to test three core hypotheses about reasoning dynamics. Our minimum viable product (MVP) implements winding number estimation, convergence metrics, persistent homology via ripser, and shape classification (settle/loop/drift). Experiments on PARARULE-Plus, synthetic counting tasks, and length-matched controls reveal that answer-token trajectories predominantly exhibit exponential convergence to fixed points, consistent with Huginn's design and independent measurements (Blayney et al., arXiv:2604.11791). No evidence of depth-dependent winding is found, refuting the naive formulation of H2. Critically, we identify \textit{steps-to-settle} as a proxy for effective computation that correlates strongly with state retention, confirmed by length-matched controls (H3 supported). Two-scale dynamics analysis (Pappone et al., NeurIPS 2025) shows that acceleration correlates with depth ($\rho = 0.997$), outperforming winding as a depth indicator. However, our reproduction of orthogonality on real FineWeb data reveals fundamentally different dynamics ($\cos \approx -0.687$, oscillatory) than reported in the original paper ($\cos \approx 0.5$--0.65, spiral), and crucially, orthogonality is not achieved without step normalization—indicating that the reported effects in Pappone et al. may be methodological artifacts rather than fundamental properties of latent dynamics. We discuss limitations, including small sample sizes and measurement methodology, and outline a roadmap for future work.
\end{abstract}

\section{Introduction}

Recent advances in test-time scaling have revived interest in recurrent-depth transformers \cite{geiping2025scaling}. Models such as Huginn employ iterative latent refinement within a fixed computational budget, raising fundamental questions about the geometry of their internal reasoning trajectories. Unlike chain-of-thought reasoning in token space, Huginn operates in a high-dimensional latent space ($\mathbb{R}^{5280}$), applying a shared recurrent block $r$ times:

\begin{equation}
    \mathbf{h}_{t+1} = \mathbf{R}(\mathbf{h}_t; \mathbf{e}),
\end{equation}

where $\mathbf{e}$ is a fixed context encoding and $\mathbf{h}_t$ represents the hidden state at recurrence step $t$. Empirically, trajectories exhibit several regimes: settling to a fixed point, looping, or drifting \cite{geiping2025scaling}. We investigate whether the geometry of these trajectories encodes computational depth or merely reflects noise.

\subsection{Theoretical Framework and Hypotheses}

The Huginn architecture initializes the latent state randomly, which "promotes convergence to a steady state independent of initialization, i.e., path independence" \cite{geiping2025scaling}. This design choice has profound implications for the model's computational capabilities. We test three hypotheses derived from the original proposal:

\begin{itemize}
    \item[H1] \textbf{Regime classification.} Each token's latent path falls into one of several regimes—settling, looping, or drifting—distinguishable via Lyapunov exponents, recurrence, and persistent homology signatures.
    \item[H2] \textbf{Loops encode reasoning depth.} When the path loops, its winding number grows with the number of reasoning steps required by the task.
    \item[H3] \textbf{Contraction prevents state tracking.} If the update is strictly contractive ($\rho(\partial_h R) < 1$), the model cannot maintain a running count or memory, requiring looping or drifting for stateful tasks.
\end{itemize}

\subsection{Related Work}

Our work builds upon several key contributions. Geiping et al. \cite{geiping2025scaling} introduced Huginn and observed latent orbits primarily on question and digit tokens. Blayney et al. \cite{blayney2026mechanistic} found non-fixed-point behavior in 0.02--2.8\% of tokens in Huginn-0125, depending on prompt conditions. Pappone et al. \cite{pappone2025two} proposed two-scale latent dynamics with spiral trajectories and orthogonality metrics. Merrill et al. \cite{merrill2024illusion} demonstrated that the "state" in state-space models is an illusion, while Grazzi et al. \cite{grazzi2024unlocking} proved that finite-precision LRNNs with positive eigenvalues cannot solve parity tasks. Yang et al. \cite{yang2025stabilizing} developed methods for stabilizing recurrent dynamics, and Tulchinskii et al. \cite{tulchinskii2025quantifying} introduced Q-K probing for logical consistency.

\section{Methods}

\subsection{Model and Protocol}

All experiments use the open Huginn-3.5B model \cite{geiping2025scaling} in inference-only mode (revision \texttt{bb6621b65e90b6a4b9b29ef88dc83866d450470c}, transformers=4.53.3). The protocol consists of:

\begin{enumerate}
    \item Forward pass through the model with hooks on \texttt{core\_block[-1]}.
    \item Extraction of hidden states at each unroll step.
    \item Computation of geometric/dynamical/topological metrics.
    \item Correlation analysis with reasoning depth and task length.
\end{enumerate}

\subsection{Metrics Implemented}

We compute nine quantities per trajectory, grouped into three categories:

\subsubsection{Geometric and Dynamical Metrics (MVP)}
\begin{itemize}
    \item \textbf{Winding number:} Accumulated angle in 2D-PCA projection, normalized to $2\pi$.
    \item \textbf{Steps-to-settle:} First step where $\|\Delta h_t\| < 0.1$, serving as a proxy for effective computation.
    \item \textbf{Shape classification:} \texttt{settle} (final step small), \texttt{loop} (return to near an earlier point), \texttt{drift} (otherwise).
    \item \textbf{Path independence:} Log-slope of $\|h_t - h_t^*\|$ for two random initializations.
    \item \textbf{Convergence rate:} Log-slope of $\|h_t - h_{\text{final}}\|$.
\end{itemize}

\subsubsection{Two-Scale Dynamics Metrics (Pappone et al., 2025)}
\begin{itemize}
    \item \textbf{Acceleration:} $a(k) = \|\Delta(k) - \Delta(k-1)\|_2$.
    \item \textbf{Orthogonality:} $\cos\theta = \frac{\Delta(k) \cdot \Delta(k-1)}{\|\Delta(k)\| \cdot \|\Delta(k-1)\|}$.
    \item \textbf{Exit step:} First step where acceleration falls below threshold for two consecutive steps.
\end{itemize}

\subsubsection{Topological Metrics}
\begin{itemize}
    \item \textbf{Persistent homology H1:} Maximum normalized persistence computed via \texttt{ripser}.
\end{itemize}

\begin{table}[h]
\centering
\caption{Summary of metrics and their implementation status}
\begin{tabular}{p{3.5cm}p{4cm}p{4cm}}
\toprule
Metric & Description & Status \\
\midrule
Lyapunov exponent & Divergence of nearby trajectories & Planned \\
Spectral radius $\rho(J)$ & Local step contraction & Planned \\
Persistent homology H1 & Topological "holes" = loops & Implemented \\
Winding number & Turns in 2D projection & Implemented \\
Steps-to-settle & Steps until $\|\Delta h\|$ decays & Implemented \\
Acceleration & Second difference of steps & Implemented \\
Orthogonality & Cosine similarity of steps & Implemented \\
\bottomrule
\end{tabular}
\end{table}

\subsection{Experimental Datasets}

We use four task families in inference-only mode:

\begin{table}[h]
\centering
\caption{Experimental datasets and configurations}
\begin{tabular}{llll}
\toprule
Task & Levels & N (trajectories) & Prompt format \\
\midrule
PARARULE-Plus (qbao775) & depths 2--5 & 10 / depth & "Question: ... Answer:" \\
Counting ($\pm1$ sum) & lengths 2--32 & 5--15 / length & "Start at 0. Add 1. Subtract 1. ..." \\
Parity (switch) & lengths 4--48 & 10 / length & "Light is off. Flip. Wait. ..." \\
Running-max & lengths 4--48 & 8 / length & "Numbers: 3 7 2 ... Largest so far?" \\
FineWeb (real texts) & -- & 1000 & -- \\
\bottomrule
\end{tabular}
\end{table}

\section{Results}

\subsection{E1: PARARULE-Plus — Winding Does Not Track Depth}

All 40 answer-token trajectories (depths 2--5, n=10 per depth) exhibit settling behavior. Correlation analysis reveals:

\begin{itemize}
    \item Winding $\sim$ depth (per-level, N=4): $\rho = +0.20$, n.s. (threshold $\approx 1.0$).
    \item Winding $\sim$ depth (per-row, N=40): $\rho = -0.037$, $p = 0.82$.
    \item Partial winding $\sim$ depth $|$ length: $\rho = -0.234$, $p = 0.15$.
    \item Steps-to-settle: $20.2 \rightarrow 22.0$ across depths (weak growth).
\end{itemize}

\textbf{Conclusion:} Multi-step symbolic logic does not induce looping; winding-depth correlation is null. Naive H2 is refuted.

\subsection{E2: Synthetic Counting — H3 Support}

Running $\pm1$ sum, length $n_{\text{ops}} \in \{2\ldots 32\}$, 5 seeds. All trajectories settle, but:

\begin{itemize}
    \item Steps-to-settle $\sim n_{\text{ops}}$: $\rho = 0.718$, $p < 1e-5$ — effective computation grows with state length.
    \item $|$winding$| \sim n_{\text{ops}}$: $\rho = 0.554$, $p = 0.002$, but partial by length yields $\rho = 0.04$, $p = 0.82$.
\end{itemize}

The winding correlation is driven by prompt length, not state tracking. This necessitates length-matched controls.

\subsection{E3: Length-Matched Control — The Killer Experiment}

\textit{Track} and \textit{local} tasks share identical body text and differ only in the question ("What is the sum?" vs. "What was the last instruction?"), ensuring identical length. With 15 seeds:

\begin{table}[h]
\centering
\caption{Length-matched control results (15 seeds)}
\begin{tabular}{lccc}
\toprule
Condition & Metric & $\rho$ & p-value \\
\midrule
\multirow{2}{*}{Track (requires accumulation)} & steps $\sim n_{\text{ops}}$ & +0.558 & $< 1e-8$ \\
 & winding $\sim n_{\text{ops}}$ & +0.324 & 0.002 \\
\multirow{2}{*}{Local (last step only)} & steps $\sim n_{\text{ops}}$ & -0.576 & $< 1e-8$ \\
 & winding $\sim n_{\text{ops}}$ & +0.054 & 0.61 (n.s.) \\
\bottomrule
\end{tabular}
\end{table}

At identical length, the sign of the correlation is determined solely by task type: state retention increases computation; pure retrieval decreases it. This constitutes strong support for H3.

\subsection{E4: Phase Map and Force-Loop}

With reduced budget (\textit{num\_steps}=16), loops appear ($n_{\text{ops}}=24$: 7/8 loop; $n_{\text{ops}}=48$: 5/8 loop). With $\textit{num\_steps} \ge 24$, all trajectories settle. The phase map (\textit{num\_steps} $\times n_{\text{ops}}$) shows that the fraction of "non-settled" trajectories grows with constrained budget and longer counts.

\textbf{Conclusion:} Loops are a symptom of insufficient computation budget, not depth per se.

\subsection{E5: Generalization to Parity}

Parity task ("light on/off"), $n_{\text{ops}} \in \{4\ldots 48\}$, 10 seeds:

\begin{itemize}
    \item Winding $\sim n_{\text{ops}}$: $\rho = 0.162$, $p = 0.22$ (n.s.).
    \item Steps-to-settle $\sim n_{\text{ops}}$: $\rho = 0.615$, $p < 1e-6$.
\end{itemize}

The same signature generalizes beyond arithmetic to other counting tasks, confirming the robustness of H3.

\subsection{Two-Scale Dynamics (Pappone et al., 2025)}

\begin{table}[h]
\centering
\caption{Two-scale dynamics on synthetic data (depths 2--5)}
\begin{tabular}{cccc}
\toprule
Depth & Mean Acceleration & Mean Orthogonality & Exits (\%) \\
\midrule
2 & $4.89 \pm 0.02$ & $-0.497 \pm 0.002$ & 50\% \\
3 & $5.84 \pm 0.01$ & $-0.498 \pm 0.001$ & 50\% \\
4 & $6.70 \pm 0.02$ & $-0.498 \pm 0.001$ & 50\% \\
5 & $7.45 \pm 0.02$ & $-0.498 \pm 0.002$ & 50\% \\
\bottomrule
\end{tabular}
\end{table}

Key findings:
\begin{itemize}
    \item \textbf{Acceleration grows with depth:} $4.89 \rightarrow 7.45$ (+52\%), correlation $\rho = 0.997$.
    \item \textbf{Orthogonality is stable:} $\sim -0.5$ across depths, indicating anti-parallel step structure.
    \item Comparison with existing metrics:
\end{itemize}

\begin{table}[h]
\centering
\caption{Comparison of depth correlations across metrics}
\begin{tabular}{lcc}
\toprule
Metric & Correlation $\rho$ & Interpretation \\
\midrule
Winding & 0.053 & Null (H2 refuted) \\
Steps-to-settle & 0.18 & Weak growth \\
Acceleration & 0.997 & Strong signal \\
Orthogonality & -0.436 & Stable, weak correlation \\
\bottomrule
\end{tabular}
\end{table}

\subsection{Reproduction of Orthogonality on Real Data: Refutation of Pappone et al.}

During reproduction of two-scale dynamics on real FineWeb text data, we obtained results that \textbf{fundamentally differ} from those reported in the original Pappone et al. paper (NeurIPS 2025, arXiv:2509.23314).

\subsubsection{Comparison with Original Paper}

\begin{table}[h]
\centering
\caption{Comparison with Pappone et al. (2025)}
\begin{tabular}{lcc}
\toprule
Aspect & Pappone et al. (2025) & Our Results \\
\midrule
Data & FineWeb (real) & FineWeb (real) \\
Dynamics type & Spiral ($\cos \approx 0.5$--0.65) & Oscillatory ($\cos \approx -0.687$) \\
Orthogonality & Achieved & Achieved \\
Step normalization & Not applied & Applied (critically important) \\
Correlation with norm & Not addressed & Resolved ($\text{corr} = -0.0436$) \\
\bottomrule
\end{tabular}
\end{table}

\subsubsection{Key Quantitative Results}

\begin{table}[h]
\centering
\caption{Metrics on real FineWeb data}
\begin{tabular}{lc}
\toprule
Parameter & Value \\
\midrule
Orthogonality ($\cos$ angle) & $-0.6870 \pm 0.0993$ \\
Angle between steps & $133.9^\circ$ \\
Acceleration & $1.8360 \pm 0.0548$ \\
Step norm & $13.7350 \pm 3.1820$ \\
Correlation with norm & $-0.0436$ (resolved) \\
Two-hit exit step & 5 (rapid stabilization) \\
\bottomrule
\end{tabular}
\end{table}

\subsubsection{Critical Observations}

During the reproduction experiment, we discovered:

\begin{enumerate}
    \item \textbf{Without step normalization, orthogonality is not stably achieved.} This implies that the orthogonality reported by Pappone et al. may be an artifact rather than a fundamental property of latent dynamics.
    \item \textbf{Correlation between orthogonality and step norm interferes with clean measurement.} Only after applying normalization do we obtain stable orthogonality.
    \item \textbf{Alternative dynamics type discovered.} Instead of spiral dynamics ($\cos \approx 0.5$--0.65), we observe oscillatory dynamics ($\cos \approx -0.687$).
\end{enumerate}

\subsubsection{Scientific Significance}

\begin{itemize}
    \item \textbf{Refutation of universality.} We demonstrate that the dynamics type (spiral vs. oscillatory) depends on measurement methodology, specifically step normalization.
    \item \textbf{Artifact identification.} Without normalization, orthogonality is not achieved, calling into question the interpretation by Pappone et al. as a fundamental property.
    \item \textbf{Methodological contribution.} Step normalization is a critical condition for clean orthogonality measurement—a factor not considered in the original paper.
\end{itemize}

\section{Discussion}

\subsection{Interpretation and Theoretical Integration}

Our results converge on a coherent picture:

1. \textbf{Convergence is by design.} Huginn's random latent initialization promotes convergence to fixed points \cite{geiping2025scaling}. All answer-token trajectories exhibit exponential contraction, consistent with Blayney et al. \cite{blayney2026mechanistic}.

2. \textbf{Counting failure is theoretically predicted.} Merrill et al. \cite{merrill2024illusion} show that "state is an illusion" in SSMs; Grazzi et al.

\cite{grazzi2024unlocking} prove that positive-eigenvalue LRNNs cannot solve parity. Huginn's contractive dynamics place it in this failure class.

3. \textbf{Steps-to-settle as effective compute proxy.} Length-matched controls conclusively show that steps-to-settle growth is driven by state retention, not prompt length. This supports H3.

4. \textbf{Refutation of Pappone et al. universality.} Our reproduction reveals that orthogonality dynamics are methodology-dependent. The original paper's conclusions may be artifacts of uncorrected measurement.

\subsection{Limitations}

\begin{itemize}
    \item \textbf{Answer-token only.} We examined only the final token; Geiping et al. observed orbits on question/digit tokens.
    \item \textbf{Small sample sizes.} Insufficient power to detect rare phenomena (0.02--2.8\% per Blayney).
    \item \textbf{Proxy metrics.} 2D-PCA winding (sign-arbitrary); steps-to-settle conflates small steps with convergence.
    \item \textbf{Single initialization seed.} Most runs use one seed; apparent signals collapsed under multi-seed testing.
    \item \textbf{Geiping regime not reproduced.} We did not use their system prompt or $r=128$.
    \item \textbf{Spectral radius $\rho(J)$ not computed.} This is the key metric for H3; we used convergence proxies.
\end{itemize}

\subsection{Future Work}

\subsubsection{Phase A: Honest Measurement}
\begin{itemize}
    \item Extend hooks to all token positions.
    \item Reproduce Geiping regime: persona prompt + $r=128$.
    \item Implement $\rho(J)$ via power-iteration + JVP \cite{yang2025stabilizing}.
    \item Reclassify regimes (settle/drift/oscillate) using $\rho(J)$.
\end{itemize}

\subsubsection{Phase B: Testing H1/H2 on Rare Events}
\begin{itemize}
    \item Power analysis to detect 0.02--2.8\% loop rates.
    \item Loop frequency scan across all tokens and 4 prompt conditions.
    \item Fix winding: global PCA (26 dims), not per-trajectory 2D.
    \item Proper TDA (population point-cloud / delay-embedding).
\end{itemize}

\subsubsection{Phase C: Mechanism}
\begin{itemize}
    \item Accuracy-gate: retain only correctly solved instances.
    \item Q-K probe \cite{tulchinskii2025quantifying} to correlate geometry with attention on PARARULE-Plus.
\end{itemize}

\section{Conclusion}

We have built a reproducible inference-only pipeline for analyzing reasoning trajectories in recurrent-depth transformers. Key findings:

\begin{itemize}
    \item Answer-token trajectories converge to fixed points (confirmed by independent metrics).
    \item The model fails synthetic counting (0\% at len $\geq 8$), consistent with theoretical predictions.
    \item No winding-depth signal detected (naive H2 refuted).
    \item \textbf{Steps-to-settle is the effective compute proxy}, correlating with state retention (H3 supported).
    \item \textbf{Two-scale dynamics:} Acceleration gives $\rho = 0.997$ vs. depth, outperforming winding.
    \item \textbf{Critical refutation:} Pappone et al.'s orthogonality results are not replicable without step normalization; we observe oscillatory dynamics ($\cos \approx -0.687$) rather than spiral ($\cos \approx 0.5$--0.65), indicating the original conclusions may be methodological artifacts.
\end{itemize}

Future work will focus on full token coverage, spectral radius computation, and power-analyzed detection of rare dynamics.

\section*{Acknowledgments}
The authors thank curator S. Barannikov and mentor N. Kurdyukov for valuable discussions and guidance. Code and data are available at \url{https://github.com/Shtirmann/Geometry-of-Reasoning-Trajectories}.

\begin{thebibliography}{9}

\bibitem{geiping2025scaling}
Geiping, J., et al. (2025). "Scaling up Test-Time Compute with Latent Reasoning: A Recurrent Depth Approach." \textit{NeurIPS 2025}. arXiv:2502.05171.

\bibitem{blayney2026mechanistic}
Blayney, E., et al. (2026). "A Mechanistic Analysis of Looped Reasoning Language Models." arXiv:2604.11791.

\bibitem{lu2025latent}
Lu, Y., et al. (2025). "Latent Chain-of-Thought? Decoding the Depth-Recurrent Transformer." arXiv:2507.02199.

\bibitem{pappone2025two}
Pappone, L., et al. (2025). "Two-Scale Latent Dynamics for Recurrent-Depth Transformers." \textit{NeurIPS 2025}. arXiv:2509.23314.

\bibitem{merrill2024illusion}
Merrill, W., Petty, J., \& Sabharwal, A. (2024). "The Illusion of State in State-Space Models." arXiv:2404.08819.

\bibitem{grazzi2024unlocking}
Grazzi, R., et al. (2024). "Unlocking State-Tracking in Linear RNNs Through Negative Eigenvalues." arXiv:2411.12537.

\bibitem{sarrof2024expressive}
Sarrof, A., Veitsman, Y., \& Hahn, M. (2024). "The Expressive Capacity of State Space Models: A Formal Language Perspective." \textit{NeurIPS 2024}. arXiv:2405.17394.

\bibitem{movahedi2026fixed}
Movahedi, S., et al. (2026). "Fixed-Point Reasoners: Stable and Adaptive Deep Looped Transformers." arXiv:2606.18206.

\bibitem{yang2025stabilizing}
Yang, H., et al. (2025). "Stabilizing Recurrent Dynamics for Test-Time Scalable Latent Reasoning in Looped Language Models." arXiv:2605.26733.

\bibitem{tulchinskii2025quantifying}
Tulchinskii, E., et al. (2025). "Quantifying Logical Consistency in Transformers via Query–Key Alignment." \textit{EMNLP 2025}. arXiv:2502.17017.

\bibitem{bai2019deep}
Bai, S., Kolter, J.Z., \& Koltun, V. (2019). "Deep Equilibrium Models." \textit{NeurIPS 2019}. arXiv:1909.01377.

\bibitem{barannikov1994framed}
Barannikov, S. (1994). "The Framed Morse Complex and its Invariants." \textit{Advances in Soviet Mathematics} 21: 93--115.

\end{thebibliography}

\newpage
\section*{Appendix: Author Contributions}

\begin{table}[h]
\centering
\caption{Contribution matrix}
\begin{tabular}{p{4cm}p{5cm}p{5cm}}
\toprule
Component & Alexandr Shianov & David Sverdlov \\
\midrule
Extraction pipeline & \checkmark & -- \\
Winding number & \checkmark & -- \\
Steps-to-settle & \checkmark & -- \\
Shape classification & \checkmark & -- \\
Persistent homology & \checkmark & -- \\
PARARULE-Plus experiments & \checkmark & -- \\
Counting experiments & \checkmark & -- \\
Length-matched controls & \checkmark & -- \\
Phase map experiments & \checkmark & -- \\
Generalization experiments & \checkmark & -- \\
Two-scale dynamics & -- & \checkmark \\
FineWeb orthogonality reproduction & -- & \checkmark \\
\hline
Primary slides (1--17) & \checkmark & -- \\
Primary slides (18--22) & -- & \checkmark \\
README documentation & \checkmark & -- \\
Two-scale README & -- & \checkmark \\
Phase A planning & \checkmark & -- \\
Phase B planning & \checkmark & -- \\
Phase C planning & -- & \checkmark \\
\bottomrule
\end{tabular}
\end{table}

\subsection*{Detailed Contributions}

\textbf{Alexandr Shianov (60\% of work):}
\begin{itemize}
    \item Developed trajectory extraction infrastructure (\texttt{extraction/}, \texttt{hook.py})
    \item Implemented all MVP metrics (\texttt{metrics/projection.py}, \texttt{winding.py}, \texttt{dynamics.py})
    \item Implemented shape classification (\texttt{shapes/gate.py}) and synthetic data generation
    \item Conducted experiments E1--E5 across PARARULE-Plus, counting, length-matched controls, phase maps, and generalization
    \item Implemented persistent homology analysis via ripser
    \item Developed correlation analysis (\texttt{analysis/correlate.py})
    \item Created main README and notebook (\texttt{01\_mvp\_h2.ipynb})
\end{itemize}

\textbf{David Sverdlov (40\% of work):}
\begin{itemize}
    \item Implemented two-scale dynamics metrics (\texttt{metrics/two\_scale.py})
    \item Conducted reproduction of orthogonality on real FineWeb data
    \item \textbf{Key discovery:} Demonstrated that without step normalization, orthogonality is not achieved; identified oscillatory dynamics ($\cos \approx -0.687$) rather than spiral ($\cos \approx 0.5$--0.65), refuting the universality of Pappone et al.'s conclusions
    \item Validated acceleration as stronger depth indicator ($\rho = 0.997$) than winding
    \item Contributed to Phase C planning (Q-K probe)
    \item Prepared slides 18--22 on two-scale dynamics and orthogonality reproduction
\end{itemize}

\end{document}
